% =============================================================================
% RL-BASED MAAS ROUTE RECOMMENDATION FOR GIUSEPPE'S MOBILITY PROBLEM
% VERSION 3: COMPREHENSIVE BEHAVIORAL & MULTI-CRITERIA FORMULATION
%   + DATA COLLECTION STRATEGY (QR tap-in/tap-out)
%   + FULL MaaS SUPER-APP (Insurance, Payment, Carpooling, Route)
%   + BEHAVIOR CHANGE PSYCHOLOGY (Nudging, Gamification)
%   + REVENUE MODEL (E-commerce, Partner Integration)
% =============================================================================
% NEXUS: A Hub for Innovation - BEST Politecnico di Torino
% March 7-8, 2026
% =============================================================================

\documentclass[11pt,a4paper]{article}

% --- Packages ---
\usepackage[utf8]{inputenc}
\usepackage[T1]{fontenc}
\usepackage[english]{babel}
\usepackage[margin=2.3cm]{geometry}
\usepackage{amsmath,amssymb,amsfonts}
\usepackage{mathtools}
\usepackage{booktabs}
\usepackage{array}
\usepackage{multirow}
\usepackage{enumitem}
\usepackage{graphicx}
\usepackage{float}
\usepackage{xcolor}
\usepackage{tcolorbox}
\tcbuselibrary{skins,breakable}
\usepackage{hyperref}
\usepackage{tikz}
\usetikzlibrary{positioning,arrows.meta,shapes,calc}
\usepackage{fancyhdr}
\usepackage{parskip}
\usepackage{textcomp}

\setlength{\headheight}{14pt}
\pagestyle{fancy}
\fancyhf{}
\fancyhead[L]{\small\textbf{RL-MaaS Framework v3 --- Full MaaS Super-App}}
\fancyhead[R]{\small NEXUS Hub -- Politecnico di Torino}
\fancyfoot[C]{\thepage}

% --- Colors ---
\definecolor{deepblue}{RGB}{25,84,123}
\definecolor{deepgreen}{RGB}{34,120,60}
\definecolor{deepred}{RGB}{180,30,30}
\definecolor{deeporange}{RGB}{220,120,20}
\definecolor{deeppurple}{RGB}{100,30,150}
\definecolor{lightblue}{RGB}{220,235,250}
\definecolor{lightgreen}{RGB}{220,250,225}
\definecolor{lightyellow}{RGB}{255,250,220}
\definecolor{lightred}{RGB}{255,230,230}
\definecolor{lightpurple}{RGB}{240,225,255}


% --- Hyperref ---
\hypersetup{colorlinks=true, linkcolor=deepblue, urlcolor=deepblue, citecolor=deepblue}

% --- Title ---
\title{
    \vspace{-1cm}
    {\normalsize\textsc{NEXUS: A Hub for Innovation -- BEST Politecnico di Torino}}\\[0.3cm]
    \rule{\textwidth}{1pt}\\[0.5cm]
    {\LARGE\textbf{MoveWise: An RL-Powered MaaS Super-App\\for Behaviorally-Aware Mobility}}\\[0.4cm]
    {\large Integrating Data Collection via QR Tap-In/Tap-Out,\\Insurance, Payment, Carpooling, and Nudge-Based\\Behavior Change for Giuseppe's Mobility Problem}\\[0.5cm]
    \rule{\textwidth}{1pt}
}

\author{
    \textbf{Team [NUMBER]}\\
    \small NEXUS 2026
}

\date{March 7--8, 2026 \quad $\vert$ \quad Version 3.0}

\begin{document}

\maketitle

% =============================================================================
% ABSTRACT
% =============================================================================
\begin{abstract}
\noindent
We present \textbf{MoveWise}, a comprehensive MaaS super-application that goes beyond route recommendation to become a full mobility ecosystem. Building on our v2 RL-based behavioral framework, v3 addresses five critical gaps: (1)~\textbf{Data collection} via a QR-code tap-in/tap-out payment system inspired by Denmark's Rejsekort/DSB, enabling us to observe mode choices, travel times, and routing patterns directly through the payment layer; (2)~\textbf{A complete MaaS service portfolio} encompassing insurance, multi-modal route planning, integrated payment, and future carpooling; (3)~\textbf{An e-commerce revenue model} where MoveWise acts as a marketplace (like Skyscanner for mobility), aggregating third-party transport services so users need only one app; (4)~\textbf{Nudge-based behavior change strategies} grounded in Thaler \& Sunstein's Nudge Theory and gamification to attract private-car users who resist mode change; (5)~\textbf{The Data Gap Problem}---addressing the fundamental challenge that car-only users generate no digital traces, and how our app solves this through service integration. The RL engine (Deep Q-Learning, HUR model, Prospect Theory, phased adoption) remains the technical core, now fed by richer behavioral data from the integrated payment and service layers.
\end{abstract}

% =============================================================================
% PROJECT FRAMING: THEME, SCOPE & CONTEXT
% =============================================================================

\begin{tcolorbox}[colback=lightpurple,colframe=deeppurple,title=\textbf{Inhabiting Time: Our Response to the NEXUS 2026 Challenge}]
\textit{``The city rushes by, but what is our pace? Moving is not just about covering a distance, it is about inhabiting time. So what happens when the transport system is not suited to our needs?''}

\vspace{4pt}
For Giuseppe, transport is already something he \textit{inhabits}: as a passenger, he studies, reads, and works during his 45-minute commute. His time is not wasted---it is lived. But this ``inhabited time'' comes at a hidden cost: pollution, dependence on a family member's schedule, and a monthly expense he finds burdensome. MoveWise does not simply shorten his trip. It \textbf{redesigns what inhabiting travel time means}: seated train journeys with Wi-Fi and power outlets for studying, gamified dashboards that transform waiting into engagement, and a phased transition that ensures he never loses the productive quality of his commute. For his errands and leisure trips---where he drives alone and time truly is wasted---MoveWise offers convenient alternatives that \textit{give that time back}.
\end{tcolorbox}

\begin{tcolorbox}[colback=lightgreen,colframe=deepgreen,title=\textbf{Target Audience, Scope \& Context}]
\begin{itemize}[noitemsep]
    \item \textbf{Target Audience:} Car-dependent periurban commuters aged 18--30 in the Turin metropolitan area, starting with university students who travel to campus from satellite towns.
    \item \textbf{Scope:} Giuseppe's \textbf{complete mobility portfolio}---not only his commute (Caselle Torinese $\rightarrow$ Orbassano, 3x/wk) but also his \textbf{shopping/errands} (3x/wk, he drives) and \textbf{leisure trips} (2x/wk, he drives). MoveWise addresses all of these.
    \item \textbf{Context of Interest:} Turin metro area, served by GTT (urban PT), SFM (regional rail), Trenitalia, and emerging micro-mobility operators (Voi, Lime). The Caselle--Orbassano corridor has no single direct PT connection, requiring intermodal solutions.
\end{itemize}
\end{tcolorbox}

% =============================================================================
% SECTION 1: THE PROBLEM --- CLEARLY STATED
% =============================================================================
\section{The Problem: Clearly Stated}
\label{sec:problem}

\subsection{Giuseppe's Situation and Preferences}

\begin{tcolorbox}[colback=lightblue,colframe=deepblue,title=\textbf{Giuseppe's Complete Profile (with all Q\&A Clarifications)}]
\begin{itemize}[noitemsep]
    \item \textbf{Who:} 23-year-old medical student
    \item \textbf{Commute:} Caselle Torinese $\rightarrow$ Orbassano, 3 times per week
    \item \textbf{Commute mode:} A \textbf{family member drives him} (he is a \textbf{passenger}, not the driver)
    \item \textbf{Commute travel time:} 45 minutes
    \item \textbf{Current monthly cost:} EUR~60/month
    \item \textbf{Desired:} 20 min shorter commute (mainly when HE drives himself)
    \item \textbf{WTP:} EUR~15/month extra for 30\% pollution reduction
    \item \textbf{Key insight:} As a passenger, he \textbf{can work/study} in the car $\Rightarrow$ his travel time is not ``wasted''
    \item \textbf{Carpooling:} NOT carpooling now, but \textbf{open to it in the future}
    \item \textbf{Values:} Views car as extremely polluting and unsafe; wants sustainability and safety
\end{itemize}
\vspace{4pt}
\textbf{His Other Trips (Often Overlooked):}
\begin{itemize}[noitemsep]
    \item \textbf{Shopping \& errands:} 3 times per week --- \textbf{he drives himself}. Short urban trips in/around Caselle Torinese.
    \item \textbf{Leisure:} 2 times per week --- \textbf{he drives himself}. Variable destinations (friends, events, recreation).
    \item \textbf{Giving lifts:} Rarely --- occasionally drives others.
\end{itemize}
\textit{These non-commute trips are where Giuseppe is the \textbf{sole polluter}. Unlike his commute (where he is a passenger and emissions are shared), his errands and leisure trips generate 100\% single-occupancy car emissions. Any complete solution must address these too.}
\end{tcolorbox}

\subsection{Considerations and Limitations}

\begin{tcolorbox}[colback=lightyellow,colframe=deeporange,title=\textbf{What We Must Consider}]
\textbf{Giuseppe's hidden advantages in his current mode:}
\begin{enumerate}[noitemsep]
    \item \textbf{Productive travel time:} He can study while being driven. Effective VOT is only $\sim$3--4 EUR~/h (vs.\ 10 EUR~/h if driving himself). Any alternative must \textit{match or exceed} this productivity.
    \item \textbf{Zero decision cost:} Someone else drives --- he doesn't plan, navigate, or transfer. Maximum convenience.
    \item \textbf{Door-to-door:} No first-mile/last-mile walking. Alternatives require access/egress time.
    \item \textbf{No technology barrier:} He doesn't use any app currently. We must onboard him.
    \item \textbf{Family relationship:} The driver is a family member --- social/emotional bond, not a transaction.
\end{enumerate}

\textbf{Limitations we face:}
\begin{enumerate}[noitemsep]
    \item \textbf{No digital data:} Car-only users leave no trace in PT/MaaS systems --- we cannot observe their travel patterns until they use our app.
    \item \textbf{Behavioral inertia:} People who always use private car are the HARDEST to convert. Habits are deeply ingrained.
    \item \textbf{Caselle--Orbassano corridor:} No single direct PT line. Requires transfers, making PT inherently less convenient.
    \item \textbf{WTP constraint:} Giuseppe is willing to pay \textit{only} EUR~15/month extra. We cannot make alternatives much more expensive.
    \item \textbf{Weather dependency:} E-scooter/cycling not viable in bad weather. Need fallback options.
\end{enumerate}
\end{tcolorbox}

\subsection{The Core Challenge}

The fundamental problem is: \textbf{how do you change the behavior of someone deeply embedded in a comfortable routine?} Giuseppe already pays EUR~60/month and finds it expensive, yet his mode is comfortable and productive. The full true costs of car travel (fuel, depreciation, insurance, externalities) far exceed the EUR~60 he perceives, but these hidden costs remain invisible to him.

Our solution has TWO parts:
\begin{enumerate}
    \item \textbf{The App (MoveWise):} A MaaS super-app that makes alternatives \textit{visible, attractive, and frictionless}.
    \item \textbf{The AI Engine (RL):} An intelligent recommendation system that learns his preferences, respects his psychology, and gradually guides him toward sustainable choices.
\end{enumerate}

% =============================================================================
% SECTION 2: DATA COLLECTION --- THE QR TAP-IN/TAP-OUT SYSTEM
% =============================================================================
\section{Creative Data Collection: QR Code Tap-In/Tap-Out}
\label{sec:data}

\subsection{The Data Gap Problem}

\begin{tcolorbox}[colback=lightred,colframe=deepred,title=\textbf{The Fundamental Data Challenge}]
When someone uses \textbf{only a private car}, we have \textbf{zero data} about their travel behavior. No tickets are purchased, no check-ins are recorded, no routes are logged. We cannot observe:
\begin{itemize}[noitemsep]
    \item Which mode they use (we can \textit{only} infer they use car)
    \item Their origin-destination pairs
    \item Their travel time
    \item Their willingness to switch
    \item Their preferences for alternative modes
\end{itemize}
This is a chicken-and-egg problem: we need data to train the RL agent, but we can only collect data if people use our app.
\end{tcolorbox}

\subsection{Solution: Inspired by Denmark's Rejsekort System}

Denmark's national travel card system, \textbf{Rejsekort} (used by DSB and all PT operators), provides a model for data-rich integrated payment:

\begin{enumerate}[noitemsep]
    \item \textbf{Tap-In (Check-In):} The user opens the MoveWise app and scans a QR code (or taps NFC) at the start of their journey. This records: timestamp, location (GPS at scan moment), selected mode.
    \item \textbf{Tap-Out (Check-Out):} The user scans again at their destination. This records: arrival timestamp, final location, mode used.
    \item \textbf{Pricing via Third-Party Providers:} MoveWise does \textbf{not} set transport prices. Fares are determined by the partner operators (GTT, Trenitalia, Voi, etc.). MoveWise aggregates and displays these prices, adding a transparent service margin. Our pricing strategy balances two goals: (a)~generating sustainable revenue for MoveWise through commissions, and (b)~keeping the total cost fair and competitive for the user compared to using each service separately.
\end{enumerate}

\textbf{What we collect from each trip:}
\begin{table}[H]
\centering
\caption{Data collected per trip via QR tap-in/tap-out}
\begin{tabular}{@{}lll@{}}
\toprule
\textbf{Data Field} & \textbf{Source} & \textbf{Use in RL} \\
\midrule
Mode used & QR scan at mode-specific reader & Revealed Preference (RP) data \\
Origin \& Destination & GPS at check-in/check-out & OD matrix construction \\
Travel time & Timestamp difference & Actual vs.\ predicted travel time \\
Price paid & Provider fare via API & Cost sensitivity estimation \\
Time of day & Timestamp & Peak/off-peak behavior \\
Weather at trip time & Cross-referenced API & Weather sensitivity \\
Trip segments & Calculated by matching sequential tap-in/tap-out events within a time window & Multi-modal trip reconstruction \\
\bottomrule
\end{tabular}
\end{table}

\subsection{How This Solves the Car-Only Data Gap}

For users who currently use \textbf{only private car} (like Giuseppe), we cannot collect this data initially. Our creative solution:

\begin{tcolorbox}[colback=lightgreen,colframe=deepgreen,title=\textbf{Strategy: Attract Car Users into the App via Additional Services}]
\begin{enumerate}[noitemsep]
    \item \textbf{Insurance service:} Offer competitive car insurance through the app (partnering with insurance companies). Car users install the app for insurance, and the app then builds their mobility profile from their interactions (searches, claims, policy details).
    \item \textbf{Parking finder:} Integrate parking availability and payment. Car users interact with the app when parking $\Rightarrow$ we observe their destination.
    \item \textbf{Fuel price comparison:} Show cheapest nearby fuel stations. Car users check the app for savings $\Rightarrow$ location data.
    \item \textbf{Carbon footprint tracker:} Show users their CO$_2$ impact per trip. Creates awareness and cognitive dissonance (they see how much they pollute).
    \item \textbf{E-commerce/Booking:} Like Skyscanner for mobility --- users search/compare/book travel options. Even if they choose car, they \textit{see} the alternatives.
\end{enumerate}
\end{tcolorbox}

\textbf{The key insight:} We don't need the user to \textit{use PT} to collect useful data. We just need them to \textit{use the app}. Once they interact with any service (insurance, parking, fuel), we start building their profile. The RL agent then uses this profile to generate increasingly relevant recommendations.

% =============================================================================
% SECTION 3: THE FULL MaaS SUPER-APP
% =============================================================================
\section{MoveWise: The Full MaaS Super-App}
\label{sec:app}

\begin{tcolorbox}[colback=lightpurple,colframe=deeppurple,title=\textbf{MoveWise Service Portfolio}]
Our app is not just a route planner. It is a \textbf{one-stop mobility ecosystem} with four core services and future expansion:
\end{tcolorbox}

\subsection{Service 1: Travel Insurance}

\begin{itemize}[noitemsep]
    \item \textbf{What:} Multi-modal travel insurance covering all modes (car, PT, e-scooter, bicycle)
    \item \textbf{How:} Partnered with insurance companies. MoveWise earns a commission per policy. Under Italian law, insurance distribution requires registration with \textbf{IVASS} (Istituto per la Vigilanza sulle Assicurazioni); MoveWise would operate as a registered digital intermediary or partner with a licensed broker, consistent with the MaaF (Mobility as a Feature) model proposed by Hensher \& Hietanen (2023).
    \item \textbf{Why it attracts car users:} Car drivers need insurance anyway. If our app offers competitive rates, they install it for insurance $\Rightarrow$ we get them into the ecosystem.
    \item \textbf{Data value:} Insurance interactions reveal user profile: vehicle type, usage frequency, risk awareness.
    \item \textbf{Nudge potential:} ``Your insurance premium would be EUR~15/month cheaper if you used PT 3 days/week'' (reward-based nudge).
\end{itemize}

\subsection{Service 2: Multi-Modal Route Planning (The AI Engine)}

\begin{itemize}[noitemsep]
    \item \textbf{What:} Intelligent route suggestion using the RL agent with behavioral awareness
    \item \textbf{How:} User enters origin-destination; the app shows ranked intermodal options with \textbf{full generalized cost} (time, money, comfort, sustainability score)
    \item \textbf{Key differentiator:} Shows not just ``fastest'' but ``best for YOU'' based on behavioral profile
    \item \textbf{Map integration:} Interactive map showing all available modes in real-time (like Google Maps + Moovit + BlaBlaCar in one)
\end{itemize}

\subsection{Service 3: Integrated Payment}

\begin{itemize}[noitemsep]
    \item \textbf{What:} One wallet for all transport modes. QR-code tap-in/tap-out for PT, automatic deduction for e-scooter/bike share, fuel payment, parking.
    \item \textbf{Models:}
    \begin{itemize}[noitemsep]
        \item \textbf{Pay-as-you-go:} Each trip charged separately (for occasional users)
        \item \textbf{Monthly subscription:} Bundled package (e.g., unlimited PT + 20 e-scooter rides + 5 carpool trips) --- like Whim/UbiGo model
    \end{itemize}
    \item \textbf{Data value:} Payment data = revealed preference. We know exactly what modes people use, how often, and how much they're willing to pay.
    \item \textbf{Rejsekort-inspired:} The QR tap-in/tap-out is directly inspired by Denmark's system. It solves the fare calculation problem for intermodal trips: you tap in when you start, tap at each transfer, and tap out at the end. The system calculates the optimal fare automatically.
\end{itemize}

\subsection{Service 4: Carpooling Platform (Future)}

\begin{itemize}[noitemsep]
    \item \textbf{What:} Connect people with similar routes to share rides. The app matches drivers and passengers automatically.
    \item \textbf{How:} Based on OD data collected from Services 1--3, the app identifies clusters of users with overlapping routes (e.g., medical students going Caselle $\rightarrow$ Orbassano).
    \item \textbf{Why future:} Giuseppe is ``not carpooling now but open to it.'' The app introduces carpooling \textit{after} trust is built (Phase 3 of adoption).
    \item \textbf{Giuseppe-specific:} His carpooling preserves the \textbf{passenger productivity benefit} --- he can still study while someone else drives. But now the emissions are halved (85 gCO$_2$/pkm vs.\ 170 for single-occupancy car).
\end{itemize}

\subsection{E-Commerce Model: ``Skyscanner for Mobility''}

\begin{tcolorbox}[colback=lightyellow,colframe=deeporange,title=\textbf{Revenue Model: How MoveWise Makes Money}]
MoveWise acts as a \textbf{marketplace/aggregator}, similar to how Skyscanner works for flights:
\begin{enumerate}[noitemsep]
    \item \textbf{Users search and book} transport services through MoveWise (bus tickets, train passes, e-scooter rentals, car insurance, parking).
    \item \textbf{Service providers pay MoveWise} a commission for each booking/referral.
    \item \textbf{Users benefit:} They don't need to install 10 different apps (GTT, Trenitalia, Voi, Lime, BlaBlaCar, etc.). One app, one account, one payment.
    \item \textbf{MoveWise earns:} Commission on bookings (5--15\%), subscription fees, insurance commissions, advertising from transport partners.
    \item \textbf{Additional revenue:} Corporate mobility packages for businesses, data analytics (anonymized) for city planning.
\end{enumerate}

\vspace{4pt}
\textbf{Revenue Estimate (Conservative):} With 10,000 active users generating an average of EUR~5/month in commissions (booking fees + insurance + subscriptions), annual platform revenue reaches \textbf{EUR~600,000/yr}. Operating costs for a lean startup (servers, 5-person team, API fees) are estimated at EUR~250,000--350,000/yr, yielding breakeven at approximately \textbf{3,000--4,000 active users}---achievable within the first year in the Turin student population alone (Politecnico + Universit\`a di Torino enroll over 100,000 students).
\end{tcolorbox}

% =============================================================================
% SECTION 3b: COMPETITORS & DISPOSAL
% =============================================================================
\subsection{Competitive Landscape}

\begin{table}[H]
\centering
\caption{MoveWise vs.\ existing mobility platforms}
\begin{tabular}{@{}p{2.2cm}p{2.5cm}p{2.5cm}p{2.5cm}p{2.5cm}@{}}
\toprule
\textbf{Capability} & \textbf{Moovit} & \textbf{Whim (Helsinki)} & \textbf{GTT Torino app} & \textbf{MoveWise} \\
\midrule
Route planning & \checkmark~(multi-city) & \checkmark & \checkmark~(GTT only) & \checkmark~(RL-personalized) \\
Integrated payment & --- & \checkmark~(subscription) & \checkmark~(GTT tickets) & \checkmark~(QR all modes) \\
Insurance & --- & --- & --- & \checkmark~(Trojan Horse) \\
Carpooling & --- & \checkmark~(taxi, not true carpool) & --- & \checkmark~(peer matching) \\
Car-user onboarding & --- & --- & --- & \checkmark~(Phase 0) \\
Behavior change & --- & --- & --- & \checkmark~(4-layer nudges) \\
Self-sustaining revenue & Ad-supported & Subscription (struggled) & Publicly funded & \checkmark~(e-commerce) \\
\bottomrule
\end{tabular}
\end{table}

\textbf{Key differentiator:} Moovit informs, Whim bundles, GTT sells tickets. None of them \textit{actively change behavior}. MoveWise is the first platform that solves the car-user data gap, deploys RL-powered personalized nudging, and funds itself through an e-commerce commission model---making it viable without the public subsidies that sank earlier MaaS pilots (Butler et al., 2021).

\subsection{Disposal \& Exit Strategy}

If MoveWise were to discontinue operations, the following assets would persist:
\begin{itemize}[noitemsep]
    \item \textbf{QR infrastructure:} Physical QR readers deployed at PT stops remain the property of the transit operator (GTT/Trenitalia) and continue to function for fare collection, independent of MoveWise.
    \item \textbf{Anonymized mobility data:} All collected OD matrices, mode-choice distributions, and travel-time data (fully anonymized under GDPR) are transferred to the \textbf{Citt\`a Metropolitana di Torino} for municipal transport planning. This makes every month of MoveWise operation a net contribution to the city's planning intelligence.
    \item \textbf{Behavioral insights:} Nudge effectiveness data (which strategies worked for which user segments) is published as an open dataset for academic research, advancing the field of MaaS adoption science.
    \item \textbf{User migration:} Users are notified 90 days in advance and guided to equivalent services (GTT app for tickets, Moovit for routing). Their QR-linked payment wallets are refunded in full.
\end{itemize}

% =============================================================================
% SECTION 4: BEHAVIOR CHANGE --- NUDGING CAR USERS
% =============================================================================
\section{Behavior Change: How to Attract Private Car Users}
\label{sec:nudge}

\subsection{The Psychology of Car Dependency}

\begin{tcolorbox}[colback=lightred,colframe=deepred,title=\textbf{Why People Resist Changing from Car}]
\textbf{Understanding the car user's mind:}
\begin{enumerate}[noitemsep]
    \item \textbf{Habit (System 1 thinking):} Car use is automatic. People don't ``decide'' to drive --- they just do it. It's like brushing teeth. (Kahneman's System 1 vs. System 2; Thaler \& Sunstein, 2008)
    \item \textbf{Loss aversion:} Switching to PT means \textit{losing} comfort, flexibility, privacy. These losses feel $2.25\times$ worse than the equivalent gains (Prospect Theory).
    \item \textbf{Status quo bias:} The default is ``keep driving.'' People stick with defaults unless actively nudged (Default Effect).
    \item \textbf{Perceived control:} In a car, you feel in control of timing, route, temperature. PT removes this control.
    \item \textbf{Social identity:} For some, a car = identity, status, freedom. Giving it up feels like losing part of yourself.
    \item \textbf{Information asymmetry:} Car users often \textbf{underestimate} the true cost of car ownership (fuel + insurance + depreciation + parking + time) and \textbf{overestimate} the cost/inconvenience of alternatives.
    \item \textbf{Lack of experience:} They've never tried PT/e-scooter/carpooling, so they fear the unknown (regret aversion).
\end{enumerate}
\end{tcolorbox}

\subsection{Our Behavior Change Toolkit}

We use a multi-layered approach based on \textbf{Nudge Theory} (Thaler \& Sunstein, 2008), \textbf{gamification}, and \textbf{social influence}:

\subsubsection{Layer 1: Nudges (Subtle Environmental Changes)}

\begin{table}[H]
\centering
\caption{Nudge strategies implemented in MoveWise}
\begin{tabular}{@{}p{3cm}p{4.5cm}p{5cm}@{}}
\toprule
\textbf{Nudge Type} & \textbf{Implementation} & \textbf{Psychological Mechanism} \\
\midrule
\textbf{Default option} & Green mode is shown \textit{first} in route results (not fastest/cheapest) & Default effect: people tend to pick the first/default option \\
\addlinespace
\textbf{Social proof} & ``87\% of students on your route use the train'' & People follow what others like them do (conformity) \\
\addlinespace
\textbf{Salience} & Show CO$_2$ impact visually (red for car, green for PT) with trees/emissions animation & Making consequences visible triggers System 2 thinking \\
\addlinespace
\textbf{Loss framing} & ``You are \textit{losing} EUR~2,100/year on hidden car costs'' instead of ``You could \textit{save} EUR~2,100'' & Loss aversion: losses hurt $2.25\times$ more (Prospect Theory) \\
\addlinespace
\textbf{Commitment device} & ``Try PT for 1 week. If you don't like it, we'll give you a free ride back'' & Reduces perceived risk; foot-in-the-door technique \\
\addlinespace
\textbf{Temporal discounting} & ``Start your green commute \textit{next Monday}'' (not today) & People agree to future commitments more easily \\
\addlinespace
\textbf{Anchoring} & Show full car cost first (EUR~510/month including all hidden costs), then PT cost (EUR~55/month) & Anchor on high number makes alternative seem cheap \\
\bottomrule
\end{tabular}
\end{table}

\subsubsection{Layer 2: Gamification (Make Sustainable Travel Fun)}

\begin{itemize}[noitemsep]
    \item \textbf{Green Points:} Earn points for every green trip. Points unlock rewards (free coffee, cinema tickets, partner discounts).
    \item \textbf{Leaderboard:} ``You saved more CO$_2$ than 73\% of users this month!'' Social competition drives behavior.
    \item \textbf{Streaks:} ``You've used PT for 5 days in a row! Keep your streak going!'' Streak psychology creates commitment.
    \item \textbf{Challenges:} Monthly challenges --- ``Try 3 different modes this month'' $\Rightarrow$ exposure to new modes reduces regret aversion.
    \item \textbf{Carbon budget:} Each user gets a monthly ``carbon budget.'' Visualization of how much they've ``spent'' creates awareness.
    \item \textbf{Badges \& milestones:} ``First train ride!'' ``10 kg CO$_2$ saved!'' ``Full multimodal week!'' --- small dopamine hits for sustainable behavior.
\end{itemize}

\subsubsection{Layer 3: Economic Incentives (Make It Financially Attractive)}

\begin{itemize}[noitemsep]
    \item \textbf{Insurance discount:} ``Use PT 3+ days/week $\Rightarrow$ 15\% lower car insurance premium'' (less driving = less accident risk).
    \item \textbf{Employer partnerships:} Companies offer MoveWise subscriptions as employee benefits (tax-deductible mobility budget).
    \item \textbf{Dynamic pricing:} Cheaper fares during off-peak hours; slightly higher during peak (spreads demand, reduces crowding).
    \item \textbf{Referral bonuses:} ``Invite a friend $\Rightarrow$ both get 1 week free PT.'' Network effects.
    \item \textbf{Cross-subsidization:} Revenue from car-related services (insurance, parking) subsidizes PT/green mode discounts.
\end{itemize}

\subsubsection{Layer 4: Education (Teach, Don't Preach)}

\begin{itemize}[noitemsep]
    \item \textbf{True Cost Calculator:} Show users the FULL cost of their car (fuel + insurance + depreciation + parking + fines + time + externalities). Most people underestimate this by 50--70\%.
    \item \textbf{``Did You Know?'' cards:} In-app educational snippets: ``Did you know? A 20-minute train ride saves 1.5 kg CO$_2$ compared to driving.''
    \item \textbf{Personalized impact report:} Monthly email: ``This month you saved X kg CO$_2$, Y EUR, and Z minutes by using MoveWise recommendations.''
    \item \textbf{Community stories:} ``Giuseppe from Caselle switched to train + e-scooter. Here's his experience.'' Peer stories are more persuasive than statistics.
\end{itemize}

\subsection{How the RL Agent Uses Nudging}

The RL agent doesn't just recommend routes --- it \textbf{chooses which nudge to deploy} based on the user's behavioral profile:

\begin{equation}
\text{Nudge}^*(i, t) = \arg\max_{\text{nudge} \in \mathcal{N}} \hat{Q}_{\text{nudge}}(s_i^t, \text{nudge}; \theta_{\text{nudge}})
\end{equation}

where $\mathcal{N}$ = \{default green, social proof, loss frame, carbon budget, streak reminder, challenge\}. A \textit{second} RL agent learns which nudge works best for which user in which context. This is \textbf{personalized nudging} --- not one-size-fits-all, but adapted to individual psychology.

For Giuseppe specifically:
\begin{itemize}[noitemsep]
    \item High $\gamma_{\text{eco}} = 0.6$ $\Rightarrow$ emphasize CO$_2$ savings (salience nudge)
    \item Segment: ``Hedonic Techy Ecologist'' $\Rightarrow$ gamification works well (points, badges)
    \item Low car-status $\gamma = 0.15$ $\Rightarrow$ don't emphasize ``give up your car'' framing
    \item High habit $H_0 = 0.7$ $\Rightarrow$ use gradual commitment device, not hard switch
\end{itemize}

% =============================================================================
% SECTION 5: GENERALIZED COST (RETAINED FROM V2)
% =============================================================================
\section{Generalized Cost Framework}
\label{sec:gc}

\textit{[This section retains all content from v2: decomposed GC with additive + non-additive costs, context-dependent VOT with productivity scores, fixed/variable cost decomposition, and Prospect Theory adjustment. See v2 document for full details.]}

Key elements:
\begin{itemize}[noitemsep]
    \item \textbf{Context-dependent VOT:} Car passenger = 3.7 EUR~/h, carpooling passenger = 5.1 EUR~/h, train (seated) = 4.4 EUR~/h, bus (seated) = 6.5 EUR~/h, driving = 10 EUR~/h
    \item \textbf{Non-additive costs:} Transfer penalties, crowding, waiting, reliability
    \item \textbf{Prospect Theory:} Loss aversion parameter $\mu = 2.25$ (Kahneman \& Tversky, 1979)
    \item \textbf{Jacobian framework:} Cross-modal cost interactions (more cars $\rightarrow$ slower buses)
\end{itemize}

% =============================================================================
% SECTION 6: BEHAVIORAL REALISM (RETAINED FROM V2)
% =============================================================================
\section{Behavioral Realism: HUR Model}
\label{sec:behavioral}

\textit{[Retained from v2: Hybrid Utility-Regret model combining utility maximization, random regret minimization (Chorus, 2008), habit persistence with exponential decay, and cultural-attitudinal factors. See v2 document for full equations.]}

Now \textbf{enhanced} with data from the QR tap-in/tap-out system:
\begin{itemize}[noitemsep]
    \item \textbf{Habit parameter calibration:} Track how quickly $H_i$ decays for real users by observing mode switching frequency in payment data.
    \item \textbf{Regret estimation:} If a user chooses train but then checks car routes 30 minutes later, this suggests regret $\Rightarrow$ high $\lambda$.
    \item \textbf{Cultural segmentation:} Cluster users by their mode choice patterns into segments (Hedonic Techy Ecologist, Neo-Luddite, Opportunist Neoclassical).
\end{itemize}

% =============================================================================
% SECTION 7: RL FORMULATION (RETAINED + ENHANCED)
% =============================================================================
\section{RL Formulation with Enhanced Data}
\label{sec:rl}

\subsection{Enhanced State Space}

The state now includes \textbf{app interaction data}:

\begin{equation}
    S_t^{\text{app}} = \left( \text{services\_used}_t,\; \text{search\_history}_t,\; \text{nudge\_response}_t,\; \text{gamification\_level}_t \right)
\end{equation}

This is added to the existing state components (network, user, MaaS, context) from v2.

\subsection{Enhanced Action Space}

The agent now recommends not just routes but also \textbf{nudges and services}:

\begin{equation}
    a_t = \left( \text{route}_t,\; \text{nudge}_t,\; \text{promotion}_t \right)
\end{equation}

\subsection{Enhanced Reward}

\begin{equation}
\boxed{
    r_t = -\left[ w_1 \cdot \widetilde{GC}^{p} + w_2 \cdot E^{p} + w_3 \cdot \Psi_{\text{behavior}}^{p} + w_4 \cdot \Phi_{\text{constraints}}^{p} \right] + \underbrace{w_5 \cdot R_{\text{revenue}}^{p}}_{\text{New: platform revenue}}
}
\end{equation}

The new term $R_{\text{revenue}}$ captures revenue generated when users book services through the app. This aligns the agent's incentives: recommend options that are good for the user, good for the environment, AND sustainable as a business.

\subsection{Constraints (Retained from v2 + 1 New)}

All 10 constraints from v2 are retained. We add:

\textbf{C11: Data Quality Constraint} --- recommendations must generate observable data:
\begin{equation}
    \mathbb{P}[\text{QR tap-in/out completed} \mid a_t] \geq \delta_{\text{data}}
\end{equation}

This ensures the agent doesn't recommend options where data collection is impossible (e.g., hitchhiking), maintaining the feedback loop.

% =============================================================================
% SECTION 8: GIUSEPPE'S SOLUTION (UPDATED)
% =============================================================================
\section{Solving Giuseppe's Problem (v3)}
\label{sec:giuseppe}

\subsection{How the App Solves Giuseppe's Preferences}

\begin{table}[H]
\centering
\caption{How MoveWise addresses each of Giuseppe's preferences}
\begin{tabular}{@{}p{3.5cm}p{3cm}p{6cm}@{}}
\toprule
\textbf{Preference} & \textbf{Challenge} & \textbf{How MoveWise Solves It} \\
\midrule
Wants shorter commute & PT requires transfers & RL agent finds optimal intermodal chains; P\&R strategy eliminates worst segments \\
\addlinespace
Values productive time & Car passenger can study & Recommend \textbf{seated train} (can study, has table/power); carpooling preserves passenger benefit \\
\addlinespace
Environmental concern & Doesn't see alternatives & \textbf{Carbon tracker} shows impact; \textbf{salience nudge} makes CO$_2$ visible; default green mode \\
\addlinespace
Safety concern & Cycling/scooter feel unsafe & Route safety scoring; avoid high-accident roads; prefer dedicated lanes; ITS integration (BSM/B2V) \\
\addlinespace
Budget constraint & WTP = EUR~15/month & Phase 2 is actually EUR~5 CHEAPER; cross-subsidized by insurance/e-commerce revenue \\
\addlinespace
Not willing to change & Strong habit ($H_0 = 0.7$) & Phased adoption (C10); nudge toolkit; gamification; gradual exposure \\
\addlinespace
No data about his travel & Car generates no PT data & Insurance \& parking services $\Rightarrow$ app installation $\Rightarrow$ interaction data $\Rightarrow$ profile building \\
\addlinespace
\textbf{Shopping/errands} (3x/wk, drives alone) & Short trips, needs to carry items & E-scooter with cargo basket for light errands; \textbf{parking finder} for car trips (captures data); fuel comparison saves money \\
\addlinespace
\textbf{Leisure trips} (2x/wk, drives alone) & Variable destinations, social context & Carpooling with friends via app (social + green); PT for central Turin destinations; gamification rewards for green leisure trips \\
\bottomrule
\end{tabular}
\end{table}

\subsection{Giuseppe's Phased Journey (Updated)}

\textbf{Phase 0 (NEW): Onboarding via Services}\\
Before Giuseppe even considers changing his commute, we get him into the app through non-transport services:
\begin{itemize}[noitemsep]
    \item Family member signs up for car insurance through MoveWise (competitive rates via partner)
    \item Giuseppe downloads the app to manage insurance documents and find parking near Orbassano
    \item The app starts building his profile from his interactions: searches, bookings, insurance claims, parking usage
    \item After 2 weeks, the app shows: ``Your commute costs you EUR~510/month in real terms (fuel + depreciation + insurance + time). Here are alternatives...''
    \item The \textbf{True Cost Calculator} creates cognitive dissonance: ``I'm paying EUR~60 but the TRUE cost is EUR~510!''
\end{itemize}

\textbf{Phase 1 (Week 1--2):} Same as v2 --- Family car to P\&R $\rightarrow$ Train. Time: 40 min, Cost: EUR~45/mo, CO$_2$: $-50\%$.

\textbf{Phase 2 (Week 3--4):} Same as v2 --- E-scooter $\rightarrow$ Train $\rightarrow$ Walk. Time: 30 min, Cost: EUR~55/mo, CO$_2$: $-75\%$.

\textbf{Phase 3 (Week 5+):} Same as v2 + Carpooling. Time: 25 min, Cost: EUR~50/mo, CO$_2$: $-80\%$.

\textbf{Payment throughout:} All trips use QR tap-in/tap-out. Giuseppe's monthly statement shows: trips taken, modes used, CO$_2$ saved, green points earned, ranking vs.\ other users.

\subsection{Beyond the Commute: Giuseppe's Other Trips}

The commute is Giuseppe's ``most important journey,'' but it accounts for only \textbf{3 of his 8+ weekly trips}. The remaining trips---where \textbf{he drives alone}---are where his personal emissions are highest and where MoveWise has the greatest marginal impact.

\begin{tcolorbox}[colback=lightyellow,colframe=deeporange,title=\textbf{MoveWise Solutions for Non-Commute Trips}]
\textbf{Shopping \& Errands (3x/wk --- he drives):}
\begin{itemize}[noitemsep]
    \item \textbf{Short trips ($<$3 km):} E-scooter with cargo attachment (Voi/Lime cargo models) or e-bike with basket. The app suggests these when the destination is nearby and the load is light.
    \item \textbf{Supermarket runs:} MoveWise's \textbf{parking finder} optimizes his car trip by pre-booking the closest spot and paying via the app---capturing his destination data even when he drives. Over time, the RL agent identifies which errands could realistically switch to micro-mobility.
    \item \textbf{Fuel price comparison:} When he does drive, the app shows the cheapest nearby station, saving money while collecting location data.
    \item \textbf{Nudge:} ``This errand is 1.2 km away. Last time you drove and spent 12 min finding parking. An e-scooter would take 6 min, door to door.''
\end{itemize}

\textbf{Leisure Trips (2x/wk --- he drives):}
\begin{itemize}[noitemsep]
    \item \textbf{Social outings:} Carpooling with friends via the app (he said he's ``open to it''). Leisure trips are the lowest-stakes context to try carpooling for the first time.
    \item \textbf{Central Turin destinations:} PT is fast and avoids parking hassle. The app highlights: ``Take the SFM train to Porta Susa---no parking needed, EUR~2.50 return.''
    \item \textbf{Gamification:} Double Green Points for sustainable leisure trips (the ``Weekend Warrior'' badge). Since leisure trips are discretionary, gamification is especially effective---there's no time pressure.
    \item \textbf{Carbon tracker:} Monthly report: ``Your leisure trips this month: 45 kg CO$_2$. If you'd used PT for the 3 central Turin trips, that would be 8 kg instead.''
\end{itemize}

\textbf{Estimated additional impact:} Switching just 2 of his 5 weekly non-commute trips to green modes saves an additional \textbf{$\sim$90 kg CO$_2$/yr}---bringing his total annual reduction to over \textbf{310 kg CO$_2$} (vs.\ 220 kg from commute alone).
\end{tcolorbox}

% =============================================================================
% SECTION 9: ENVIRONMENTAL IMPACT (RETAINED FROM V2)
% =============================================================================
\section{Environmental Impact}
\label{sec:impact}

\textit{[Retained from v2: per-user savings (220 kgCO$_2$/yr, EUR~4,100/yr generalized cost), system-level impact (10,000 users: $-$2,200 tonnes CO$_2$/yr, EUR~1.6M/yr societal benefit). Full externality framework: CO$_2$ + noise + accidents + congestion + health.]}

% =============================================================================
% SECTION 10: DIFFERENCES FROM V2
% =============================================================================
\section{What Changed: v2 $\rightarrow$ v3}
\label{sec:differences}

\begin{table}[H]
\centering
\caption{v2 vs.\ v3 comparison}
\begin{tabular}{@{}p{3cm}p{5cm}p{5cm}@{}}
\toprule
\textbf{Aspect} & \textbf{v2} & \textbf{v3 (This)} \\
\midrule
Data collection & Assumed available & QR tap-in/tap-out system (Rejsekort-inspired) \\
\addlinespace
Car-user data gap & Not addressed & Solved via service integration (insurance, parking, fuel) \\
\addlinespace
App scope & Route recommendation only & Full MaaS super-app (insurance, payment, route, carpool) \\
\addlinespace
Revenue model & Not considered & E-commerce aggregator (Skyscanner model) \\
\addlinespace
Behavior change & Passive (HUR model) & Active (nudges, gamification, education, economic incentives) \\
\addlinespace
Onboarding & Assumed user installs app & Phase 0: attract via non-transport services \\
\addlinespace
RL enhancement & 4-component reward & 5-component reward (+ revenue); nudge selection agent \\
\addlinespace
Problem statement & Implicit & Explicitly stated with considerations and limitations \\
\addlinespace
Non-commute trips & Not addressed & Shopping, errands, leisure trips explicitly covered with dedicated solutions \\
\addlinespace
Competitors \& disposal & Not addressed & Competitive analysis (Moovit, Whim, GTT) and exit strategy included \\
\addlinespace
Regulatory awareness & Not addressed & IVASS insurance regulation and revenue breakeven validated \\
\bottomrule
\end{tabular}
\end{table}

% =============================================================================
% SECTION 11: CONCLUSION
% =============================================================================
\section{Conclusion}

MoveWise v3 transforms from a route recommendation system into a \textbf{complete MaaS super-app} that addresses the full lifecycle of behavior change:

\begin{enumerate}
    \item \textbf{Attract:} Get car users into the app through non-transport services (insurance, parking, fuel comparison).
    \item \textbf{Educate:} Show them the true cost of car dependency and the benefits of alternatives (True Cost Calculator, CO$_2$ tracker).
    \item \textbf{Nudge:} Use behavioral psychology (default green, social proof, loss framing, gamification) to make sustainable choices the path of least resistance.
    \item \textbf{Recommend:} The RL agent suggests personalized, phased, psychology-aware routes that users will \textit{actually follow}.
    \item \textbf{Learn:} QR tap-in/tap-out collects rich behavioral data that improves the RL agent over time.
    \item \textbf{Sustain:} Revenue from the e-commerce model (commissions, subscriptions, insurance) makes the platform financially viable without requiring public subsidies.
\end{enumerate}

\begin{tcolorbox}[colback=lightyellow,colframe=deeporange,title=\textbf{The One-Sentence Pitch (v3)}]
\centering
\textit{``MoveWise is a MaaS super-app that uses Rejsekort-inspired QR ticketing, nudge psychology, and Reinforcement Learning to transform private car users into sustainable commuters --- not by telling them what to do, but by making the green choice the easy, fun, and financially smart choice.''}
\end{tcolorbox}

% =============================================================================
% REFERENCES
% =============================================================================
\section*{References}

\begin{itemize}[leftmargin=*, label={}, nosep]
    \item Chorus, C.G., Arentze, T.A., \& Timmermans, H.J.P. (2008). A Random Regret-Minimization Model of Travel Choice. \textit{Transportation Research Part B}, 42(1), 1--18.
    \item Kahneman, D., \& Tversky, A. (1979). Prospect Theory: An Analysis of Decision under Risk. \textit{Econometrica}, 47(2), 263--291.
    \item Mnih, V. et al.\ (2015). Human-Level Control through Deep Reinforcement Learning. \textit{Nature}, 518, 529--533.
    \item Thaler, R.H., \& Sunstein, C.R. (2008). \textit{Nudge: Improving Decisions about Health, Wealth, and Happiness}. Yale University Press.
    \item Wang, L., Duan, P., Lyu, C., \& Ma, Z. (2025). RL-based Sequential Route Recommendation for SO Traffic Assignment. \textit{arXiv:2505.20889v1}.
    \item Rejsekort A/S (2026). Denmark's National Travel Card System. \url{https://www.rejsekort.dk/}
    \item DSB (2026). Danish State Railways Digital Ticketing. \url{https://www.dsb.dk/}
    \item Wardman, M. (2004). Public Transport Values of Time. \textit{Transport Policy}, 11(4), 363--377.
    \item Sochor, J., Str\"omberg, H., \& Karlsson, M. (2014). Travelers' Motives for Adopting a New, Innovative Travel Service: Insights from the UbiGo Field Operational Test. \textit{21st ITS World Congress}, Detroit.
    \item UITP (2019). \textit{Policy Brief: Mobility as a Service}. Brussels.
    \item Butler, L., Yigitcanlar, T., \& Paz, A. (2021). Barriers and Risks of MaaS Adoption. \textit{Cities}, 109, 103036.
    \item Carbonara, N. et al.\ (2024). Business Model Innovation in MaaS. \textit{J.\ Cleaner Production}, 464, 142744.
    \item Jittrapirom, P. et al.\ (2017). MaaS: A Critical Review. \textit{Urban Planning}, 2(2), 13--25.
    \item Pronello, C. (2025--2026). Course Lectures: Transport Innovation for Sustainable, Inclusive and Smart Mobility. Politecnico di Torino.
    \item EEA (2024). European Environment Agency: Transport Emission Factors.
    \item ERSO (2024). European Road Safety Observatory: Cyclist Safety Report.
    \item Hensher, D.A., \& Hietanen, S. (2023). Mobility as a Feature (MaaF): Rethinking the Role of Integration. \textit{Transport Reviews}, 43(3), 365--381.
    \item Hensher, D.A., \& Nelson, J.D. (2025). MaaF: Insurance Rewards and Emissions Reduction. \textit{Transportation Research Part A}, 181, 103912.
\end{itemize}

\end{document}
